\documentclass[11pt]{article}
\usepackage[margin=1in]{geometry}
\usepackage{hyperref}
\usepackage{enumitem}
\hypersetup{colorlinks=true,urlcolor=blue}
\pagestyle{empty}

\begin{document}

\begin{flushright}
Shuaidong Gao \\
Chongqing Institute of Foreign Studies \\
Qijiang, Chongqing 401420, China \\
\texttt{gaoshuaidong@stu.cqifs.edu.cn} \\
ORCID: 0009-0004-5641-3581 \\
\today
\end{flushright}

\bigskip

\noindent Dear KDD 2027 Program Chairs and Area Chairs,

\medskip

I am submitting {\bf causalscale: Scaling Differentiable Causal Discovery to Genome-Wide Resolution} to the Datasets \& Benchmarks Track.

I want to be upfront: I am an independent researcher at a teaching-focused university in western China. I work alone, on a laptop with 8 GB of VRAM. Every result in this paper---seven engines, a 33-cancer pan-cancer scan, DepMap genome-wide validation, and cross-domain experiments spanning genomics, finance, and climate---ran on a single RTX 5060.

That is, in a way, the point. \texttt{causalscale} exists because the tools I needed did not: no single package could go from DAGMA at 30 variables through Causal Transformer at 200--500 to LowRankGNN at 17,787 genes, all through one \texttt{fit()} call. The same function that recovers ARID1A--MTOR tissue-specific directionality in breast cancer also discovers sector structure in S\&P 500 equities and the west-to-east ENSO propagation in NOAA climate reanalysis. I built it so that researchers who cannot afford compute clusters could still do causal discovery at scale.

The paper includes:

\begin{itemize}[nosep]
\item A unified Python package integrating 7 engines (DAGMA, ClusterAware, Causal Transformer, LowRankGNN, MultiScale, MultiModal, Ensemble) under automatic dimension-based selection
\item PCMCI time-series engine (\texttt{tigramite} wrapper) for lagged causal discovery, verified on synthetic VAR(1) and NOAA NCEP monthly temperature
\item Two-stage pipeline (DAG engine as seed finder $\to$ LowRankGNN as expander) for genome-scale discovery beyond the DAG ceiling
\item Top-K sparse correlation training (O($K$) memory, 14$\times$ reduction), inspired by CorALS (Becker et al., Nature Comp.\ Sci.\ 2023)
\item Synthetic benchmarks against NOTEARS and DAGMA (matched seeds, matched hardware, 5 seeds per configuration)
\item STRING/TRRUST external validation at 88.7\% precision (574/647 edges) on DepMap genome-wide data
\item 33-cancer pan-cancer scan recovering tissue-specific ARID1A--MTOR directionality (Spearman $\rho=-0.720$, $p=2.3{\times}10^{-6}$)
\item Three-domain cross-validation: genomics (TCGA/DepMap), finance (S\&P 500, 76\% same-sector), climate (NOAA NCEP, ENSO west-to-east propagation)
\item Ensemble voting weight ablation (125 configurations) and engine boundary crossover analysis ($d=190$--$210$, $d=490$--$510$)
\item One-command reproduction (\texttt{python run\_all.py --verify}), 92\% test coverage
\item Pre-trained causal graphs on HuggingFace Hub for researchers without GPU access
\end{itemize}

I have done my best to be honest about limitations: LowRankGNN relaxes the DAG constraint at genome scale, ASCEND validation uses a small sample, and three engines are based on my own SSRN preprints currently under review. The paper acknowledges the $O(d^2)$ correlation bottleneck at $d=10^8$ and reports where v3.2.0 top-K sparse training and the two-stage pipeline partially mitigate it. Ensemble weights and boundary selection were systematically ablated rather than left as heuristics. References include recent KDD work on scalable differentiable CD (Ma et al., KDD~2024), dynamic causal discovery (Wang \& Song, KDD~2025), and time-series benchmarking (Ferdous \& Gani, KDD~2025).

\medskip

\noindent {\bf Suggested Reviewers}

\smallskip

I respectfully suggest the following five reviewers, all of whom have published directly in differentiable causal discovery, structure learning, or the computational biology data used in this work:

\begin{enumerate}[nosep]
\item {\bf Xun Zheng} --- Carnegie Mellon University \\
      Author of NOTEARS (NeurIPS 2018), the foundational method on which \texttt{causalscale}'s ClusterAware engine is built. \\
      \texttt{xunzheng@andrew.cmu.edu} (verified via Google Scholar)

\item {\bf Bryon Aragam} --- University of Chicago, Booth School of Business \\
      Co-author of DAGMA (NeurIPS 2022) and NOTEARS. His work on differentiable structure learning directly underpins Engine 7. \\
      \texttt{bryonaragam@uchicago.edu} (publicly listed on UChicago faculty page)

\item {\bf Ignavier Ng} --- Carnegie Mellon University \\
      Author of GOLEM (NeurIPS 2020) and the differentiable causal discovery survey (2023). Deep expertise in the theoretical foundations of score-based DAG learning. \\
      \texttt{ignavierng@cmu.edu} (verified via Google Scholar)

\item {\bf Philippe Brouillard} --- Mila, Universit\'{e} de Montr\'{e}al \\
      Author of ``Differentiable Causal Discovery from Interventional Data'' (NeurIPS 2020). Expertise in extending differentiable methods to interventional settings---a natural next step for \texttt{causalscale}. \\
      \texttt{philippe.brouillard@umontreal.ca} (verified via Google Scholar)

\item {\bf Aviad Tsherniak} --- Broad Institute of MIT and Harvard \\
      Lead scientist on the Cancer Dependency Map (DepMap) project. \texttt{causalscale} uses DepMap CRISPR data for its genome-scale validation; his perspective on whether the biological results are meaningful would be invaluable. \\
      \texttt{aviad@broadinstitute.org} (publicly listed via Broad Institute)
\end{enumerate}

\bigskip

\noindent Thank you for your time and for considering this submission. KDD's commitment to open, reproducible research made it the natural venue for a paper whose core contribution is making causal discovery accessible to researchers who, like me, work at the edge of institutional infrastructure. I believe this paper meets the Datasets \& Benchmarks Track's bar for software impact, reproducibility, and community value.

\medskip

\noindent Sincerely,

\medskip

\noindent Shuaidong Gao

\end{document}
